\documentclass[11pt]{article}

\usepackage[a4paper,margin=1in]{geometry}
\usepackage[T1]{fontenc}
\usepackage[utf8]{inputenc}
\usepackage{lmodern}
\usepackage{microtype}
\usepackage{booktabs}
\usepackage{longtable}
\usepackage{array}
\usepackage{xcolor}
\usepackage{hyperref}
\usepackage{xurl}
\usepackage{enumitem}
\usepackage{amsmath}
\usepackage{listings}

\hypersetup{
  colorlinks=true,
  linkcolor=blue,
  urlcolor=blue,
  citecolor=blue
}

\definecolor{codebg}{RGB}{246,248,250}
\lstset{
  basicstyle=\ttfamily\small,
  backgroundcolor=\color{codebg},
  breaklines=true,
  frame=single,
  rulecolor=\color{gray!35},
  columns=fullflexible
}

\setlist[itemize]{topsep=2pt,itemsep=2pt}
\setlist[enumerate]{topsep=2pt,itemsep=2pt}
\renewcommand{\arraystretch}{1.15}
\Urlmuskip=0mu plus 2mu

\title{\textbf{Multimodal Radar Image Iceberg Classification}\\
\large SDSC8007 Deep Learning Mini Project}
\author{
Yang Lin, 72540214 \\
Duan Yixuan, 72542270
}
\date{\today}

\begin{document}
\maketitle

\begin{center}
\begin{tabular}{lp{0.62\linewidth}}
\toprule
Course & SDSC8007 Deep Learning \\
Competition & Statoil/C-CORE Iceberg Classifier Challenge \\
Selected topic & Intermediate Task 3: Multimodal Radar Image Iceberg Classification \\
\bottomrule
\end{tabular}
\end{center}

\section*{Abstract}

This project reproduces and improves the Kaggle Statoil/C-CORE Iceberg Classifier Challenge, a binary classification task that predicts whether a target in satellite radar imagery is an iceberg or a ship. Each sample contains two Synthetic Aperture Radar (SAR) image bands and an incidence-angle metadata feature, so the task is naturally a multimodal classification problem.

We built a complete modeling pipeline covering data analysis, leakage-aware validation, baseline construction, CNN training, model improvement, final result analysis, AI-usage disclosure, and reproducibility packaging. Starting from a shallow engineered-feature baseline, we progressively added SAR image CNNs, pseudo-label filtering, a pretrained FiLM ResNet34, incidence-angle statistical features, and strict second-stage stacking.

The final model is a strict angle-aware stacking ensemble. The first stage trains multiple CNN image models with 5-fold cross-validation. The second stage combines their out-of-fold predictions with incidence-angle statistical features using LightGBM and Logistic Regression. To reduce validation leakage, when predicting each validation fold, the true labels of that fold are excluded from label-based incidence-angle statistics.

The final Kaggle late submission achieved a private score of \textbf{0.08098} and a public score of \textbf{0.09987}. The corresponding local strict 5-fold cross-validation log loss is \textbf{0.088556}, below the Top-5-level reference value 0.08883 used in our project comparison.

\section*{Author Contributions}

Both team members made equal contributions to this mini-project. Yang Lin and Duan Yixuan jointly contributed to task selection, literature and public-solution review, model design, experiment execution, result analysis, presentation preparation, and final report writing. Therefore, we report the contribution breakdown as:

\begin{center}
\begin{tabular}{lc}
\toprule
Member & Contribution \\
\midrule
Yang Lin & 50\% \\
Duan Yixuan & 50\% \\
\bottomrule
\end{tabular}
\end{center}

\section{Task Description \& Metric}

The selected competition is the Statoil/C-CORE Iceberg Classifier Challenge. The goal is to classify objects in satellite SAR images as either iceberg or ship. Each sample contains two radar image bands and one incidence-angle metadata feature.

\begin{center}
\begin{tabular}{ll}
\toprule
Field & Description \\
\midrule
\texttt{band\_1} & First SAR image band, flattened from a $75 \times 75$ image \\
\texttt{band\_2} & Second SAR image band, also flattened from a $75 \times 75$ image \\
\texttt{inc\_angle} & Radar incidence angle, provided as numerical metadata \\
\texttt{is\_iceberg} & Binary target label; 1 indicates iceberg and 0 indicates ship \\
\bottomrule
\end{tabular}
\end{center}

This is a binary classification task with multimodal inputs. It is more challenging than standard natural-image classification because SAR images represent radar backscatter rather than RGB color and texture. The task matches the course's Intermediate Task 3, which emphasizes non-natural images, multimodal feature fusion, and noise or robustness analysis.

The official evaluation metric is binary log loss:

\[
\mathrm{LogLoss}=-\frac{1}{N}\sum_{i=1}^{N}\left[y_i\log(p_i)+(1-y_i)\log(1-p_i)\right],
\]

where $y_i$ is the true label and $p_i$ is the predicted probability that the sample is an iceberg. Lower log loss is better. Since log loss heavily penalizes overconfident wrong predictions, the model must output calibrated probabilities rather than only correct class labels.

Kaggle late submission is available for this competition. We therefore report the official Kaggle public and private scores for the final submission, while also reporting a strict 5-fold CV score to document the validation protocol and reproducibility.

\section{Data Analysis \& Challenges}

The training set contains 1,604 labeled samples, and the test set contains 8,424 samples. Each sample has two SAR bands of size $75 \times 75$. There are 133 missing incidence-angle values in the training set and no missing incidence-angle values in the test set.

Our cache-building script reports:

\begin{lstlisting}
train images: (1604, 2, 75, 75), labels: (1604,)
test images:  (8424, 2, 75, 75)
train missing incidence angles: 133
test missing incidence angles: 0
test real-like inc_angle decimals <= 4: 3424
test synthetic-like inc_angle decimals > 4: 5000
\end{lstlisting}

The main challenges are as follows.

\begin{itemize}
  \item \textbf{Non-natural SAR imagery.} SAR images measure radar backscatter rather than visible-light appearance. They contain speckle-like noise and have different statistics from standard RGB images.
  \item \textbf{Small training set.} Only 1,604 labeled samples are available, so deep CNNs can overfit easily without careful validation, regularization, data augmentation, and ensembling.
  \item \textbf{Multimodal structure.} The target depends on both image morphology and incidence angle. The incidence angle affects radar reflection and also contains strong dataset-level statistical patterns.
  \item \textbf{Machine-generated test rows.} Public discussions reported that the test set contains generated rows. We preserved the raw \texttt{inc\_angle} string and computed its decimal precision. In our audit, 3,424 test rows have \texttt{angle\_decimals <= 4}, while 5,000 rows have \texttt{angle\_decimals > 4}. This information is used as a filtering signal for pseudo-labeling and angle-neighborhood statistics.
  \item \textbf{Public/private instability.} The public leaderboard uses only part of the test labels. Public and private scores can differ noticeably, so we treat the private score as the main Kaggle evidence and use strict CV as a reproducibility check.
\end{itemize}

\section{Validation Strategy}

We use a two-layer validation strategy: first-stage validation for CNN image models and second-stage validation for angle-aware stacking.

\subsection{First-stage CNN validation}

All first-stage CNN models are trained using stratified 5-fold cross-validation.

\begin{center}
\begin{tabular}{ll}
\toprule
Setting & Value \\
\midrule
Number of folds & 5 \\
Shuffle & True \\
Random seed & 2026 \\
Split strategy & Stratified by \texttt{is\_iceberg} \\
Metric & Binary log loss \\
\bottomrule
\end{tabular}
\end{center}

For each fold, image normalization and incidence-angle normalization are fitted only on the training portion of that fold, then applied to the validation portion and test set. This prevents validation statistics from entering the training process.

\subsection{Strict second-stage stacking validation}

The second stage uses first-stage out-of-fold predictions and incidence-angle-derived statistical features. The central leakage-control rule is:

\begin{quote}
When predicting a validation fold, the true labels of that validation fold are not used to construct label-based incidence-angle statistics for that fold.
\end{quote}

For example, when fold 0 is used as validation, statistics such as the iceberg rate of samples with the same or nearby incidence angle are computed only from folds 1--4. The model can use unlabeled \texttt{inc\_angle} values from train and test as input metadata, but it cannot use validation-fold labels when constructing label-based features.

\subsection{Leakage-control summary}

\begin{center}
\begin{tabular}{p{0.43\linewidth}p{0.47\linewidth}}
\toprule
Potential risk & Control measure \\
\midrule
Image normalization fitted with validation statistics & Fit normalization only on the training split of each fold \\
Incidence-angle scaler fitted with validation statistics & Fit the angle scaler only on the training split \\
Label-based angle statistics leaking validation labels & Exclude the current validation fold's labels in strict stacking \\
Pseudo-label noise from synthetic-like test rows & Use only high-confidence pseudo labels and filter by raw \texttt{inc\_angle} decimal precision \\
Overconfident probabilities under log loss & Clip final probabilities and optimize blend weights on out-of-fold predictions \\
\bottomrule
\end{tabular}
\end{center}

Base C uses pseudo-labeling. It does not use any test-set true labels, but it does use unlabeled test inputs and teacher-model predictions. We therefore disclose it as a semi-supervised and transductive training strategy.

\section{Baseline Model}

The baseline model is designed to be simple, reproducible, and interpretable. We compute image statistics for each SAR channel, including mean, standard deviation, minimum, maximum, and percentile values. These features are combined with \texttt{inc\_angle} and a missing-angle indicator, then trained using Logistic Regression and ExtraTrees. Their out-of-fold predictions are blended.

\begin{center}
\begin{tabular}{p{0.24\linewidth}p{0.56\linewidth}r}
\toprule
Model & Validation setup & Log loss \\
\midrule
Logistic Regression + ExtraTrees blend & 5-fold CV & 0.3403 \\
\bottomrule
\end{tabular}
\end{center}

This baseline verifies that data loading, label alignment, fold splitting, and log-loss computation are correct. It also shows that global image statistics alone are not sufficient for Top-5-level performance.

\section{Model Improvements}

\subsection{SAR image channel construction}

The original input contains two SAR image bands. We construct a 4-channel dB-space input:

\begin{lstlisting}
channel 1 = band_1
channel 2 = band_2
channel 3 = (band_1 + band_2) / 2
channel 4 = band_1 - band_2
\end{lstlisting}

The first two channels preserve the raw SAR bands, the third channel represents average radar intensity, and the fourth channel captures polarization contrast. Channel-wise normalization is computed within each fold using only the training split.

\subsection{First-stage image models}

We train four first-stage model families.

\begin{center}
\begin{tabular}{p{0.16\linewidth}p{0.46\linewidth}p{0.27\linewidth}}
\toprule
Model & Method & Purpose \\
\midrule
Base A & Compact residual-style CNN with 4-channel dB input and an angle auxiliary branch & Stable image baseline \\
Base B & VGG-style CNN with batch normalization and dropout & Architectural diversity \\
Base C & CNN trained with filtered high-confidence pseudo labels from Base A & Semi-supervised signal expansion \\
Base D & ImageNet-pretrained ResNet34 with FiLM angle modulation & Pretrained representation and physical angle conditioning \\
\bottomrule
\end{tabular}
\end{center}

All models use rotations, flips, small Gaussian noise, early stopping, and test-time augmentation. The purpose of the first stage is not to make one CNN solve the task alone, but to generate diverse and stable out-of-fold predictions for the second-stage model.

\subsection{Pretrained model usage}

Base D uses the official PyTorch ResNet34 ImageNet checkpoint:

\begin{lstlisting}
https://download.pytorch.org/models/resnet34-b627a593.pth
\end{lstlisting}

ResNet34 has approximately 21--22 million parameters, which is far below the course limit of 0.5B parameters. Since SAR inputs are not RGB images, the first convolution is modified from 3 input channels to 4 input channels. The first three channels are initialized from the pretrained RGB weights, while the fourth channel is initialized by averaging the RGB weights. We then add a FiLM module that uses \texttt{inc\_angle} to generate feature-wise scale and shift parameters for the CNN feature map.

\subsection{Incidence-angle statistical features and stacking}

Public top-solution discussions show that \texttt{inc\_angle} contains strong group and neighborhood patterns. We did not copy public notebooks. Instead, we reimplemented this idea with a stricter validation rule that excludes validation-fold labels.

The second-stage feature set includes:

\begin{itemize}
  \item predictions from the four first-stage CNN models;
  \item mean, median, minimum, maximum, range, and standard deviation of CNN predictions;
  \item raw \texttt{inc\_angle} value;
  \item missing-angle indicator;
  \item raw JSON decimal precision of \texttt{inc\_angle};
  \item prediction statistics for identical or rounded incidence angles;
  \item KNN neighborhood statistics in incidence-angle space;
  \item radius-neighborhood statistics in incidence-angle space;
  \item label-based angle statistics computed under the strict fold rule.
\end{itemize}

The second-stage ensemble contains:

\begin{lstlisting}
4 LightGBM stacking models, seeds = 2026, 2027, 2028, 2029
4 Logistic Regression stacking models, seeds = 2026, 2027, 2028, 2029
\end{lstlisting}

Their out-of-fold predictions are combined using non-negative weight optimization. The final model is an 8-model strict stacking blend.

\subsection{Improvement summary}

\begin{center}
\begin{tabular}{p{0.24\linewidth}p{0.56\linewidth}r}
\toprule
Stage & Main method & Log loss \\
\midrule
Baseline & Image statistics + \texttt{inc\_angle} + traditional models & 0.3403 \\
CNN/image ensemble & Multiple CNN image models & 0.1788 \\
Strict angle-aware stacking & CNN predictions + strict \texttt{inc\_angle} statistics + LightGBM/LogReg blend & 0.088556 \\
\bottomrule
\end{tabular}
\end{center}

The CNN ensemble significantly improves over the shallow baseline. The largest final gain comes from explicitly modeling incidence-angle structure after obtaining stable image-based predictions.

\section{Final Results}

\subsection{Local strict CV result}

The final strict 8-model blend achieves:

\begin{lstlisting}
Strict 5-fold CV log loss = 0.088556
\end{lstlisting}

This is below the Top-5-level reference value 0.08883 used in our project comparison.

\subsection{Kaggle late-submission result}

The final main submission file is:

\begin{lstlisting}
submission_blend_8models_20260503_212025.csv
\end{lstlisting}

Kaggle late-submission scores:

\begin{center}
\begin{tabular}{lrr}
\toprule
Submission & Private Score & Public Score \\
\midrule
\texttt{submission\_blend\_8models\_20260503\_212025.csv} & \textbf{0.08098} & 0.09987 \\
\bottomrule
\end{tabular}
\end{center}

Because public/private instability is substantial in this competition, the private score is the main Kaggle evidence. We do not claim an official current Top-5 leaderboard rank unless the leaderboard page explicitly shows it. Instead, we state that the achieved log loss reaches the course Top-5-level performance target.

\section{AI Tool Usage Statement}

AI tools were used as assistants, but not to directly generate and submit a complete solution. The usage was limited to:

\begin{enumerate}
  \item \textbf{Debugging and error diagnosis.} AI helped diagnose Python, PyTorch, CUDA/GPU, dependency, server-path, and shell-script errors.
  \item \textbf{Technical explanation and idea discussion.} AI helped explain public discussion ideas such as incidence-angle statistics, stacking, pseudo-labeling, and FiLM, and helped analyze their applicability and leakage risks.
  \item \textbf{Code review and reproducibility checks.} AI helped inspect the cross-validation protocol, pseudo-label filtering, second-stage stacking, README, and one-command reproduction script.
  \item \textbf{Report organization and language polishing.} AI helped organize the report and polish the wording, while all reported experimental results, Kaggle scores, method choices, and limitations were checked by the group.
\end{enumerate}

We did not directly copy or lightly modify public Kaggle notebooks, and we did not submit a complete AI-generated implementation. Public discussions were used only as references for general modeling ideas. The group reimplemented the data pipeline, model training, cross-validation, strict stacking, final blending, and reproduction scripts for this project.

Core designs independently confirmed by the group include:

\begin{itemize}
  \item binary log-loss evaluation and 5-fold CV setup;
  \item SAR channel construction and fold-safe normalization;
  \item first-stage CNN training configuration;
  \item strict angle-aware stacking with validation-fold label exclusion;
  \item filtered pseudo-labeling and its semi-supervised/transductive nature;
  \item source, parameter count, and structure of pretrained ResNet34 + FiLM;
  \item final Kaggle late submission and reproduction workflow.
\end{itemize}

\section{Discussion \& Reflection}

The experiments show that training a deeper CNN alone is not sufficient for this task. The shallow baseline obtains 0.3403 log loss. CNN image models improve the result to around 0.1788, confirming that SAR image morphology is useful. The final major improvement comes from combining image predictions with strict incidence-angle statistical features, reducing local CV log loss to 0.088556 and Kaggle private score to 0.08098.

Effective components include fold-safe CNN training and test-time augmentation; multiple first-stage image models rather than a single CNN; explicit incidence-angle group and neighborhood features; second-stage validation that excludes current validation-fold labels; and LightGBM/Logistic Regression blending for stable probability estimates.

Limitations remain. The incidence-angle pattern is dataset-specific and may not generalize to a new radar acquisition process. Base C uses test inputs and teacher predictions for pseudo-labeling, so it should be disclosed as semi-supervised/transductive training. Final blend weights are optimized on out-of-fold predictions and may contain mild model-selection bias. CNN retraining can also vary slightly across CUDA/PyTorch environments even with fixed seeds.

Future work could include a no-pseudo-label ablation, systematic SAR speckle-noise filtering, nested validation for blend-weight selection, and deeper public/private score analysis.

\section{Reproducibility}

\subsection{Code package}

The reproduction package is:

\begin{lstlisting}
SDSC8007_Iceberg_Classifier_Repro.zip
\end{lstlisting}

The public GitHub repository for the same code package is:

\begin{lstlisting}
https://github.com/mayday4yl/SDSC8007-Iceberg-Classifier
\end{lstlisting}

It does not include Kaggle raw data, old experiment outputs, or model weights. To reproduce, place the Kaggle files in:

\begin{lstlisting}
data/processed/train.json
data/processed/test.json
data/processed/sample_submission.csv
\end{lstlisting}

\subsection{Environment}

The final reproduction environment was:

\begin{center}
\begin{tabular}{ll}
\toprule
Item & Configuration \\
\midrule
GPU & NVIDIA GeForce RTX 4090 24GB \\
CUDA & 12.8 \\
Driver & 570.124.06 \\
Python & 3.12 \\
PyTorch & 2.6.0a0+ecf3bae40a.nv25.01 \\
Seed & 2026 \\
\bottomrule
\end{tabular}
\end{center}

Python dependencies are listed in \texttt{requirements.txt}. PyTorch and torchvision should be installed using a CUDA-compatible build for the target machine.

\subsection{One-command reproduction}

After placing the Kaggle data files under \texttt{data/processed/}, run:

\begin{lstlisting}[language=bash]
cd /root/SDSC8007_Iceberg_Classifier_Repro
python -m pip install -U pip
python -m pip install -r requirements.txt
PYTHON_BIN=python bash scripts/run_from_scratch.sh
\end{lstlisting}

The script creates a timestamped output directory:

\begin{lstlisting}
artifacts/from_scratch_<timestamp>/
  run.log
  predictions/
  reports/
  models/
\end{lstlisting}

Key outputs include:

\begin{lstlisting}
predictions/submission_blend_8models_*.csv
predictions/oof_blend_8models_*.csv
reports/metrics_blend_8models_*.json
run.log
\end{lstlisting}

The Kaggle submission CSV is a generated output. The file submitted for late submission was named \texttt{submission\_blend\_8models\_20260503\_212025.csv}; the filename itself does not affect Kaggle scoring, but we keep the original name in the report so it matches the Kaggle submission record.

The final CV score can be recomputed from the generated OOF file using the verification script in the \texttt{scripts} directory.

\section*{References}

\begin{enumerate}
  \item Kaggle. Statoil/C-CORE Iceberg Classifier Challenge. \url{https://www.kaggle.com/competitions/statoil-iceberg-classifier-challenge}
  \item David Austin. 1st Place Solution Overview. Kaggle Discussion 48241. \url{https://www.kaggle.com/competitions/statoil-iceberg-classifier-challenge/discussion/48241}
  \item beluga. 2nd Place Solution Overview. Kaggle Discussion 48294. \url{https://www.kaggle.com/competitions/statoil-iceberg-classifier-challenge/discussion/48294}
  \item Evgeny Nekrasov. 3rd-Place Solution Overview. Kaggle Discussion 48207. \url{https://www.kaggle.com/competitions/statoil-iceberg-classifier-challenge/discussion/48207}
  \item Andrii Sydorchuk. 4th Place Approach. Kaggle Discussion 48163. \url{https://www.kaggle.com/competitions/statoil-iceberg-classifier-challenge/discussion/48163}
  \item Azat Akhtyamov. 6th Place Solution. Kaggle Discussion 48193. \url{https://www.kaggle.com/competitions/statoil-iceberg-classifier-challenge/discussion/48193}
  \item PyTorch. ResNet34 ImageNet pretrained weights. \url{https://download.pytorch.org/models/resnet34-b627a593.pth}
\end{enumerate}

\end{document}
